\documentclass[../main.tex]{subfiles}
\begin{document}

% --------------------------------------------------
\section{math.sty}
% --------------------------------------------------

\texttt{math.sty}では, 数式入力を簡単にするためのパッケージとマクロをまとめている. 
\texttt{amssymb}, \texttt{mathrsfs}, \texttt{mathtools}を読み込み, よく使うアルファベット, 作用素, 括弧, アクセントを短く書けるようにしている. 

\subsection{ブラックボード・カリグラフィック・フラクタル・スクリプト}

\texttt{mathbb}, \texttt{mathcal}, \texttt{mathfrak}, \texttt{mathscr}のアルファベットは, 全て簡潔に書くためのマクロを用意している. 
\begin{itemize}
	\item \verb|\mathbb{Z}| : \verb|\bbZ|
	\item \verb|\mathcal{C}| : \verb|\calC|
	\item \verb|\mathfrak{p}| : \verb|\frakp|
	\item \verb|\mathscr{F}| : \verb|\scrF|
\end{itemize}

\begin{ex}{}{}
	\begin{align}
		\bbZ, \quad \calC, \quad \frakp, \quad \scrF.
	\end{align}
\end{ex}

\subsection{数学作用素}

プリアンブルに\verb|\MyMathOperators{}|と記入し, 要素として\texttt{Ker}等の数学作用素を入れると, \verb|\mathop{\mathrm{Ker}}|を定義したのと同じように使える. 
例えば
\begin{center}
	\verb|\MyMathOperators{Ker,Hom,Gal,Spec,rank}|
\end{center}
とプリアンブルに書いたとすると
\begin{align}
	\Ker f, \quad \Hom(E_1, E_2), \quad \Gal(L/K), \quad \Spec \bbC[x], \quad \rank E(\bbQ)
\end{align}
のように書くことができる.

\begin{rem}{}{}
	数学作用素の間に
	\begin{center}
		\verb|\MyMathOperators{Ker, Hom, Gal, Spec, rank}|
	\end{center}
	のような半角スペースを入れるとエラーが生じる. 
	カンマの直後にはスペースを入れずに書く. 
\end{rem}

\subsection{括弧・絶対値・集合マクロ}

絶対値や括弧を始めとする様々な記号は以下のように書く.

\begin{itemize}
	\item \verb|\abs{x}| : 絶対値
	\item \verb|\norm{x}| : ノルム
	\item \verb|\rbra{x}| : 丸括弧
	\item \verb|\cbra{x}| : 波括弧
	\item \verb|\sbra{x}| : 角括弧
	\item \verb|\abra{x}| : 内積または角括弧
	\item \verb|\floor{x}| : 床関数
	\item \verb|\ceil{x}| : 天井関数
	\item \verb|\set{x\in X}{P(x)}| : 集合
\end{itemize}

\begin{rem}{}{}
	\texttt{mathtools}の\verb|\DeclarePairedDelimiter|を使って定義しているので, マクロ名の後に\texttt{*}を付けると中身に合わせて括弧の大きさが自動調整される. 
	\begin{align}
		\rbra{\frac{b}{a}}, \quad \rbra*{\frac{b}{a}}.
	\end{align}
\end{rem}

\subsection{アクセント}

長い数式に対して通常のアクセント\verb|\bar{}|, \verb|\tilde{}|, \verb|\hat{}|を付けると, 線や記号の幅が不足することがある. 
そこで\texttt{math.sty}では, それぞれ\verb|\overline{}|, \verb|\widetilde{}|, \verb|\widehat{}|が実行されるように再定義している. 

\begin{ex}{}{}
	\begin{align}
		\bar{a+bi}=a-bi.
	\end{align}
\end{ex}

\subsection{数式番号}

数式番号は, 引用した数式にのみ付くようにしている. 
\texttt{mathtools}の\texttt{showonlyrefs}を使った設定である. 

\begin{ex}{}{}
	式
	\begin{align}
		e^{i\pi}+1=0 \label{Euler}
	\end{align}
	は有名な等式である. \zcref{Euler}はオイラーの等式と呼ばれる.
\end{ex}

\subsection{記号の調整}

\LaTeX 標準の\verb|\epsilon|は\verb|\varepsilon|に再定義している. 
そのため, 本文中で\verb|\epsilon|と書くと, 常に$\epsilon$が出力される. 

また, Tate--Shafarevich群などで使う記号として\verb|\Sha|を定義している. 
\begin{ex}{}{}
	\begin{align}
		\Sha(E/K)
	\end{align}
\end{ex}

\end{document}
